\section{Kokkuvõte} \label{kokkuvõte}

Käesoleva bakalaureusetöö eesmärk oli kavandada ja realiseerida suurel keelemudelil põhinev prototüüp, mis annab eestikeelsele informaatika lõputöö tekstile automaatset tagasisidet struktuuri, akadeemilise stiili, terminoloogia järjepidevuse ja viitamisvajaduse kohta, ning hinnata empiiriliselt selle väljundi kvaliteeti. Töös määratleti Tartu Ülikooli arvutiteaduse instituudi nõuete~\cite{ut_juhend_2024} põhjal operatsionaalselt neli tagasisidekategooriat, kavandati kaks promptimisstrateegiat (üldine ja struktureeritud) ning loodi veebipõhine prototüüp, mis võimaldab kasutada kahte tippmudelit (Claude 4.7 Opus ja GPT-5.5) eestikeelse teksti analüüsiks. Empiirilises osas viidi läbi nii kvalitatiivne pilootuuring kui ka täismahuline kvantitatiivne hindamine 80 päris API-päringu põhjal 20 katkendist koosneval sünteetilisel testkogul.

\textbf{Empiirilised tulemused} (alapeatükk~\ref{hindamine:synteetiline}) annavad uurimisküsimustele 1--3 järgmised vastused:
\begin{enumerate}
    \item \textbf{Mudeli tagasiside iseloom (UK1).} Mõlemad mudelid annavad auditeeritavat eestikeelset tagasisidet kõigis neljas kategoorias, kuid kalduvad eelistama saagist täpsusele (saagis 0,77--0,93, täpsus 0,27--0,47), mistõttu süsteem sobib kandidaatide leidmise abivahendiks, kuid selle tagasiside vajab kasutaja kriitilist kontrolli.
    \item \textbf{Kategooriate raskus (UK2).} Akadeemilise stiili kategoorias on F\textsubscript{1} kõrgeim (0,52--0,74), struktuuri puhul madalaim (Claude'i üldise prompti puhul 0,11), terminoloogia ja viitamisvajadus jäävad vahepeale (0,34--0,60).
    \item \textbf{Promptitüübi mõju (UK3).} Struktureeritud prompt parandab makro-F\textsubscript{1}-skoori järjepidevalt: Claude'il +0,05 ja GPT-5.5-l +0,19. Parim kombinatsioon on GPT-5.5 koos struktureeritud promptiga (F\textsubscript{1} = 0,56), mis kinnitab töö keskset hüpoteesi struktureeritud promptimise eelisest.
\end{enumerate}

\textbf{Töö panus.} Töö peamine panus on (1) struktureeritud LLM-prompti disain eestikeelse informaatikateksti tagasisidestamiseks, mille eelist üldise prompti ees näidatakse empiiriliselt; (2) kategooriapõhine empiiriline võrdlus Claude 4.7 Opuse ja GPT-5.5 vahel eestikeelses kontekstis, kus avaldatud kvantitatiivseid võrdlusi pole varem dokumenteeritud; (3) avatud prototüüp koos reprodutseeritava demo-režiimi ja täielikult dokumenteeritud hindamisprotsessiga, mida saab uuesti rakendada päris üliõpilastöödel.

\textbf{Piirangud ja edasine töö.} Peamised piirangud (alapeatükk~\ref{hindamine:piirangud}) on testkogu sünteetilisus ja kuldstandardi tuginemine ainult ühe annoteerija (autori) hinnangule. Jätkukava (alapeatükk~\ref{hindamine:edasine}) näeb esmalt ette sama hindamise kordamist päris üliõpilastööde katkenditel, seejärel sõltumatu teise annoteerija kaasamist koos Coheni $\kappa$~\cite{cohen_kappa_1960} arvutamisega ning kategooriapõhise prompti edasist täpsustamist. Kokkuvõttes näitab töö empiiriliselt, et tipptasemel LLM-id on jõudmas tasemele, kus nende abil saab anda eestikeelsele akadeemilisele tehnilisele tekstile praktikas kasutuskõlblikku tagasisidet. Süsteemi roll on seejuures selgelt täiendav, mitte asendav: lõputöö autoriks jääb üliõpilane ning juhendaja sisuline tagasiside on endiselt asendamatu.
